% !TEX root = ../jlt.tex
\section{Evaluation}
\label{app:fjlt:evaluation}

Evaluation is small-step and call-by-value on closed terms, $e \longrightarrow e'$, with values and evaluation contexts
\[
\begin{array}{rcl}
v &::=& c \mid f \mid O\texttt{<}\overline{\tau}\texttt{>\{}\overline{v}\texttt{\}} \mid \texttt{std::pair\{}v\texttt{, }v\texttt{\}} \mid \iota_1\,v \mid \iota_2\,v \mid \texttt{std::optional\{}v\texttt{\}} \mid \texttt{std::nullopt} \\
  &\mid& \texttt{std::make\_shared<const }\tau\texttt{>(}v\texttt{)} \\
E &::=& [\,] \mid f\texttt{<}\overline{\sigma}\texttt{>(}\overline{v}\texttt{, }E\texttt{, }\overline{e}\texttt{)} \mid E\texttt{(}\overline{e}\texttt{)} \mid v\texttt{(}\overline{v}\texttt{, }E\texttt{, }\overline{e}\texttt{)} \mid E \oplus e \mid v \oplus E \mid \tau\texttt{(}E\texttt{)} \mid E \texttt{ ? } e \texttt{ : } e \\
  &\mid& \iota_i\,E \mid \texttt{case(}E\texttt{, }e\texttt{, }e\texttt{)} \mid \texttt{case(}v\texttt{, }E\texttt{, }e\texttt{)} \mid \texttt{case(}v\texttt{, }v\texttt{, }E\texttt{)} \mid \ldots
\end{array}
\]
where the elided contexts are the remaining argument positions of pairs, optionals, boxes, projections, constructions, folds and the two branches of a case, each evaluated left to right.
It is \Cpp{}'s constant evaluation restricted to the grammar, extended by \textsc{S-Unbox} at run time, and Figure~\ref{fig:fjlt-eval} gives the rules; $e \Downarrow v$ abbreviates $e \longrightarrow^{*} v$.
The context grammar picks one redex per non-value term, so the relation is deterministic.

\begin{figure}[htbp]
\footnotesize
\[
\begin{array}{ll}
\textsc{(S-Call)} & f\texttt{<}\overline{\sigma}\texttt{>(}\overline{v}\texttt{)} \longrightarrow e[\overline{\sigma}/\overline{\alpha}][\overline{v}/\overline{x}] \\
 & \qquad \text{where } f\texttt{<}\overline{\sigma}\texttt{>} \text{ resolves to a declaration with parameters } \overline{x} \text{ and body } e \\[3pt]
\textsc{(S-Op)} & O\texttt{<}\overline{\tau}\texttt{>\{}\overline{v_y}\texttt{\}(}\overline{v}\texttt{)} \longrightarrow e[\overline{\tau}/\overline{\alpha}][\overline{\sigma}/\overline{\beta}][\overline{v_y}/\overline{y}][\overline{v}/\overline{x}] \quad \text{for a declared operation with body } e\text{, } \overline{\sigma} \text{ deduced} \\[3pt]
\textsc{(S-Std)} & O\texttt{<}\overline{\tau}\texttt{>\{\}(}\overline{c}\texttt{)} \longrightarrow c \quad \text{for a standard function object, } c \text{ the value the standard assigns} \\[3pt]
\textsc{(S-Prim)} & c_1 \oplus c_2 \longrightarrow c \quad \text{the base type's operation, total by the choice of } \oplus \\[3pt]
\textsc{(S-Conv)} & \tau\texttt{(}c\texttt{)} \longrightarrow c' \quad \text{the conversion between distinct carriers the standard or the declaration assigns} \qquad \tau\texttt{(}v\texttt{)} \longrightarrow v \quad \text{at } v\text{'s own type, any value} \\[3pt]
\textsc{(S-If)} & \texttt{true ? } e_1 \texttt{ : } e_2 \longrightarrow e_1 \qquad \texttt{false ? } e_1 \texttt{ : } e_2 \longrightarrow e_2 \\[3pt]
\textsc{(S-Proj)} & \texttt{std::pair\{}v_1\texttt{, }v_2\texttt{\}.first} \longrightarrow v_1 \qquad \texttt{std::pair\{}v_1\texttt{, }v_2\texttt{\}.second} \longrightarrow v_2 \qquad O\texttt{<}\overline{\tau}\texttt{>\{}\overline{v}\texttt{\}.}y_i \longrightarrow v_i \\[3pt]
\textsc{(S-Unbox)} & \texttt{*make\_shared<const }\tau\texttt{>(}v\texttt{)} \longrightarrow v \\[3pt]
\textsc{(S-Case)} & \texttt{case(}\iota_1\,v\texttt{, }v_f\texttt{, }v_g\texttt{)} \longrightarrow v_f\texttt{(}v\texttt{)} \qquad \texttt{case(}\iota_2\,v\texttt{, }v_f\texttt{, }v_g\texttt{)} \longrightarrow v_g\texttt{(}v\texttt{)} \\[3pt]
\textsc{(S-Range)} & \texttt{size(}r\texttt{)} \longrightarrow |r| \qquad \texttt{head(}r\texttt{)} \longrightarrow \mathit{head}\,r \text{ if } |r| > 0 \text{, } \mathit{default}_\tau \text{ if } |r| = 0 \qquad \texttt{next(}r\texttt{)} \longrightarrow \mathit{next}\,r \text{ if } |r| > 0 \text{, } r \text{ if } |r| = 0 \\[3pt]
\textsc{(S-Fold)} & \texttt{fold\_left(}r\texttt{, }v_a\texttt{, }v_f\texttt{)} \longrightarrow \texttt{fold\_left(}\mathit{next}\,r\texttt{, }v_f\texttt{(}v_a\texttt{, }\mathit{head}\,r\texttt{), }v_f\texttt{)} \quad \text{if } |r| > 0 \\
 & \texttt{fold\_left(}r\texttt{, }v_a\texttt{, }v_f\texttt{)} \longrightarrow v_a \quad \text{if } |r| = 0 \\[3pt]
\textsc{(S-Cata)} & \texttt{cata(}v_{\mathit{alg}}\texttt{, }n\texttt{)} \longrightarrow v_{\mathit{alg}}\texttt{(}n\texttt{, cata(}v_{\mathit{alg}}\texttt{, }l_1\texttt{), }\ldots\texttt{, cata(}v_{\mathit{alg}}\texttt{, }l_k\texttt{))} \\
 & \qquad \text{for a node } n \text{ with legs } l_1 \ldots l_k \\[3pt]
\textsc{(S-Ctx)} & E[e] \longrightarrow E[e'] \quad \text{if } e \longrightarrow e' \\
\end{array}
\]
\normalsize
\caption{Evaluation.
\textsc{S-Call} covers specialisations and generic functions alike, since resolution has already chosen the body; for a generic function the explicit type arguments are substituted first.
\textsc{S-Range} is total: on the empty range the head is $\mathit{default}_\tau$ and the tail the range itself, values \textsc{S-Fold} never reaches, since it runs $|r|$ times.
The default value of a type is given by its structure, not by \cppinline{T{}}: a carrier's own default, a pair of defaults, the empty optional, the first injection of a default, a fieldless function object, and for a box a box of its content's default, so that no term of the kernel holds a null box.
\textsc{S-Fold} consumes the head of the range, so a fold takes as many steps as the range has elements.
\textsc{S-Cata} unfolds one node, passing the node itself beside its legs' folds; a leaf has no legs and the algebra is applied to it alone.
\textsc{S-Prim} and \textsc{S-Std} are total because the admitted operations are, which is what Lemma~\ref{lem:progress} needs.
\textsc{S-Unbox} is evaluation at run time only: \cppinline{std::shared_ptr} is not \cppinline{constexpr} in \Cpp{23}, so a constant evaluator rejects a term that builds a box, and a law or a \cppinline{static_assert} premise contains no \cppinline{make_shared}, which is the condition under which every FJlt program is a Jlt program.}
\label{fig:fjlt-eval}
\end{figure}

